\chapter{Implementation}
\label{chapter:implementation}

% =============================================================================
\begin{comment}
\section{Vehicle Configuration and Baseline Mission}


\subsection{First-Stage Engine Model}
\label{ssec:engine_model}

Thrust and specific impulse switch at the staging events, and the second stage flies its
vacuum values, less the pressure deficit of its nozzle in the cases of Chapter~\ref{ssec:case_definitions}
that carry it. The first stage is flown with one of two models of the
ambient-pressure dependence of Chapter~\ref{ssec:thrust}, summarized in
Table~\ref{tab:engine_models}.

The pressure-dependent model, flown in every case but one, evaluates
Eq.~\eqref{eq:thrust_pressure_alt} with the nozzle exit area of Table~\ref{tab:stage_params}.
As in Chapter~\ref{ssec:thrust}, its specific impulse is linear in the ambient pressure, and is
scaled between its published endpoints,
\begin{equation}
  I_{sp}(h) = I_{sp,\mathrm{vac}}
     - \frac{p_a(h)}{p_0}\,\bigl(I_{sp,\mathrm{vac}} - I_{sp,SL}\bigr) \ ,
  \label{eq:isp_pressure}
\end{equation}
with the mass flow following as $\dot m_p = T/(g_0 I_{sp})$. A constant mass flow requires
$T_{SL}/T_{vac} = I_{sp,SL}/I_{sp,vac}$; the published figures give $0.925$ and $0.910$. They are
kept as published, so the mass flow falls by \qty{1.6}{\percent} over the burn. This is the only model
under which the pressure loss of Chapter~\ref{sec:performance} is defined.

The sea-level model holds the lift-off values for the whole burn. It is the most conservative
constant choice: with the same propellant and the highest mass flow, it ends the first-stage
burn earliest and delivers the least impulse.

\begin{table}[htbp]
  \centering
  \caption{The two first-stage engine models.}
  \label{tab:engine_models}
  \begin{tabular}{l l l}
    \toprule
     & \textbf{Pressure-dependent (baseline)} & \textbf{Sea level, constant} \\
    \midrule
    Thrust           & \qty{7.61}{\mega\newton} at lift-off to \qty{8.23}{\mega\newton} in vacuum & \qty{7.61}{\mega\newton} throughout \\
    Specific impulse & \qtyrange{283}{311}{\s}                             & \qty{283}{\s} throughout \\
    Mass flow        & \qtyrange{2741}{2698}{\kg\per\s} (\qty{-1.6}{\percent})       & \qty{2741}{\kg\per\s} throughout \\
    First-stage burn & \qty{146.0}{\s}                                    & \qty{144.4}{\s} \\
    Pressure loss    & defined                   & not defined \\
    \bottomrule
  \end{tabular}
\end{table}

\subsection{Payload Fairing}
\label{ssec:fairing_jettison}

The \qty{1900}{\kg} payload fairing~\cite{NASA_Falcon9v12_Datasheet_2018} is booked within the first-stage empty mass of
Table~\ref{tab:stage_params} rather than added to it; stage separation drops only the rest of
that mass. It is jettisoned, and the atmosphere taken as left, as the vehicle climbs through
\qty{65}{\km}~\cite{kluever2018space}. Of the criteria of Chapter~\ref{ssec:fairing}, the altitude
is preferred because it
is the same for every trajectory, whereas a load- or heating-based threshold would move the
jettison with the variables being optimized. The criterion governs the jettison alone: drag
is modeled throughout the first-stage flight, and the exo-atmospheric laws engage at
second-stage ignition (Chapter~\ref{sec:numerical_methods}).

The jettison is an instantaneous drop in mass at the altitude crossing, handled as an arc
boundary, like staging (Chapter~\ref{ssec:integration})~\cite{BettsSurvey1998}. If the first
stage burns out below \qty{65}{\km}, the drop is applied at the next arc boundary, stage separation
or second-stage ignition, so the fairing is never carried under thrust above that altitude.

\end{comment}

% =============================================================================
\begin{comment}


\section{Software Implementation}
\label{sec:software}

\subsection{Architecture}
\label{ssec:architecture}

The code is organized as a layered dependency graph. At the bottom sit leaf data
and models that depend on nothing else: the physical constants and vehicle
specifications, and the environment models (gravity, atmosphere, Earth rotation).
Above them the configuration layer exposes every user-facing switch. The guidance
laws are also leaves --- pure functions that map a state to a steering command ---
so that dependencies always run from the simulation core towards the guidance
modules and never the reverse, which keeps each law testable in isolation. At the
center is the physics hub \texttt{rocket\_ascent.py}, which imports all guidance
and auxiliary modules and integrates the equations of motion; the solvers and
optimizers wrap the hub, and the plotting layer consumes plain result arrays,
decoupled from the solvers.

The solver layer is not flat, and the way it is stacked matters.
\texttt{pso\_coast\_solver.py} owns the second-stage vacuum dynamics and the
per-trajectory guidance state described in the next subsection, and two further solvers
are built \emph{on} it rather than beside it: \texttt{direct\_pso\_solver.py} imports its
integrator and guidance machinery and contributes only a different arc layout and objective,
and \texttt{segmented\_guidance\_solver.py} does likewise while drawing its waypoints from
\texttt{segment\_reference.py}, which builds the indirect reference trajectory once and caches
it to disk, keyed on the target orbit, the vehicle, the environment and the budget of its
solve. The cache is the only
piece of state the tool carries between runs. Figure~\ref{fig:dependency_tree} shows the
resulting layering.

\begin{figure}[h!]
\centering
\includegraphics[width=0.95\linewidth]{Figures/module_dependency_graph.png}
\caption{Module dependency graph.}
\label{fig:dependency_tree}
\end{figure}

Every architecture of Chapter~\ref{sec:optimization}, and the segmented mode, wraps the
same inner engine in its outer optimizer; only the design vector, the termination
condition and the objective differ. Figure~\ref{fig:flowchart} shows the loop for the
apogee-circularization architecture (A), and Figure~\ref{fig:arc_structures} the arc structures of the
others (B to E).

The numerical machinery is provided by established libraries rather than
reimplemented: the Dormand--Prince integrator, the root finders, and the grid search come
from \texttt{SciPy}~\cite{virtanen2020scipy}, and the particle swarm optimizer from
\texttt{PyGMO}~\cite{biscani2020pagmo}.

\subsection{Guidance Dispatch and Trajectory State}
\label{ssec:dispatch}

Section~\ref{sec:laws_implemented} stated that every guidance law is a stateless
function and that the memory it needs is held by the caller. Where that memory
lives is the one structural feature of the implementation that a reader would not
guess, and it explains a behavior reported later, so it is set out here.

The nine laws are driven through \emph{two} independent dispatchers. The
apogee-circularization architecture steers from inside the equations-of-motion routine itself
(\texttt{rocket\_ascent.rocket\_dynamics}), which selects the active law and keeps
its working variables --- the explicit-guidance steering constants, the epoch at
which they were last recomputed, the coefficient-freeze flags, the constant pitch
rate --- in module-level variables, reset in a block at the start of each run. The
three population-based architectures and the segmented mode instead steer from
\texttt{pso\_coast\_solver}, which carries exactly the same variables in a
\texttt{GuidanceState} object created fresh for each candidate trajectory and
re-initialized at every arc boundary.

The duplication is deliberate rather than accidental. A swarm evaluates many
thousands of trajectories inside a single process, so any guidance variable cached
at module level survives from one candidate to the next and quietly corrupts the
evaluation that follows; binding that state to an object which is discarded with
the particle removes the failure mode by construction. The same hazard exists in
the shared first-stage code, where a running engine-performance ramp has to be
explicitly reset at the start of every trajectory for precisely this reason.

Where the closed-loop laws recompute their steering coefficients matters for the same reason.
Recomputed inside the equations-of-motion routine, they would also be updated at the trial
points at which the adaptive integrator evaluates the derivatives before accepting or
rejecting a step --- instants the trajectory may never reach. In every case of
Chapter~\ref{chapter:results_and_discussion} they are recomputed once per guidance cycle, on
the state the integrator has accepted.

Two consequences follow, and both are relevant when reading the results. First, all
three population-based architectures share a single first-stage routine, which
always applies the instantaneous pitch-over of Chapter~\ref{ssec:kick}; the
triangular kick profile is therefore exercised only by the apogee-circularization
architecture, whatever the configuration requests. Second, a law can behave
differently under different architectures even though the law itself is unchanged.
The documented case is the constant-pitch-rate law, whose ramp starts from the
vertical in the apogee-circularization path but from the current flight-path angle in the
population-based path, so that the commanded angle of attack is continuous at
ignition; the two produce visibly different trajectories from the same nominal
guidance selection. Differences of this kind are properties of the dispatch, not of
the steering logic, and are noted wherever they affect a comparison.

A third consequence proved more serious, and is recorded here because the way it arose is
instructive rather than incidental. The equations of motion are implemented in two layers: a
small routine holding the dynamics of Eqs.~\eqref{eq:planar_eom_a}--\eqref{eq:planar_eom_e}, which is
by construction the \emph{non-rotating} model, and above it the dispatcher, which adds the
rotating-frame corrections of Eqs.~\eqref{eq:sim_dv_pseudo}--\eqref{eq:sim_dgamma_pseudo}
before returning. The population-based solvers cannot call the dispatcher for their upper
stage --- it carries the guidance state, event side-effects and history logging described
above, none of which can be executed thousands of times inside a swarm --- so they call the
inner routine directly. In doing so they inherited the dynamics but not the corrections
applied above them, and their upper stage was integrated in a rotating frame without the
terms that frame requires. Nothing failed: the omission produced no error, and the optimizer
simply converged to a slightly different control.

The corrections are now applied inside the shared upper-stage routine as well, so that no
trajectory is integrated under one force model before staging and another after it; the
optimal-control architecture carries them too (Chapter~\ref{ssec:indirect_pmp_impl}). The general
point applies to any further development of the tool: a term added to the dispatcher rather
than to the shared inner routine will be missing, silently, from every architecture that
bypasses the dispatcher.

    
\end{comment}

\begin{comment}
    


\subsection{Configuration}
\label{ssec:configuration}

All user-facing behavior is controlled from
\texttt{Input\_File/simulation\_parameters.py}, without editing the solver code.
Table~\ref{tab:config} groups the principal switches by function; each mission is
fully specified by a choice within these groups plus the vehicle and
mission-geometry parameters of Section~\ref{sec:vehicle}.

\begin{table}[htbp]
  \centering
  \caption{Principal configuration switches, grouped by function.}
  \label{tab:config}
  \begin{tabular}{l p{0.7\textwidth}}
    \toprule
    \textbf{Group} & \textbf{Representative parameters} \\
    \midrule
    Guidance selection      & \texttt{GUIDANCE\_MODE} (nine laws),
                              \texttt{TGO\_ESTIMATOR}, \texttt{PEG\_CONVERGENCE\_MODE},
                              \texttt{CPR\_THETA\_DOT\_MODE},
                              \texttt{GUIDANCE\_REFRESH\_MODE} \\
    Mission architecture    & \texttt{COAST\_METHOD}
                              (\texttt{apogee\_check}/\texttt{pso\_coast}/\texttt{direct}/%
                              \texttt{reference\_track}),
                              \texttt{DIRECT\_LAW\_TERMINATED\_CUTOFF},
                              \texttt{DIRECT\_OPTIMIZER} \\
    Segmented guidance      & \texttt{MULTI\_GUIDANCE\_ENABLED},
                              \texttt{GUIDANCE\_SEGMENTS},
                              \texttt{MULTI\_GUIDANCE\_OPTIMIZE\_ALTITUDES},
                              \texttt{MULTI\_GUIDANCE\_ALT\_LB}/\texttt{\_UB},
                              \texttt{SEGMENTED\_LAW\_TERMINATED\_ARCS} \\
    Reference cache         & \texttt{PMP\_REFERENCE\_USE\_CACHE},
                              \texttt{PMP\_REFERENCE\_FORCE\_RERUN},
                              \texttt{PMP\_REFERENCE\_PSO\_PARTICLES}/%
                              \texttt{\_GENERATIONS}
                              (fidelity of the cached reference) \\
    Earth-rotation / frames & \texttt{ENABLE\_EARTH\_ROTATION},
                              \texttt{INCLUDE\_PSEUDO\_FORCES},
                              \texttt{AZIMUTH\_INCLINATION\_MODE} \\
    Atmosphere / aero       & \texttt{INCLUDE\_DRAG}, \texttt{INCLUDE\_LIFT},
                              \texttt{ATMOSPHERE\_EXIT\_METHOD}
                              (altitude / dynamic pressure / aerothermal flux),
                              \texttt{FAIRING\_JETTISON\_MODE} \\
    Kick profile            & \texttt{KICK\_PROFILE\_MODE} (triangular / instantaneous;
                              triangular honored only under \texttt{apogee\_check}) \\
    Engine modes            & \texttt{ISP\_1\_MODE}, \texttt{THRUST\_1\_MODE}
                              (sea-level / vacuum / average / linear / pressure;
                              set the two alike) \\
    Optimizer settings      & \texttt{RUN\_FAST}; four independent PSO blocks
                              (indirect, coast, direct, segmented), each with
                              $\omega$, $c_1$, $c_2$, $v_{\max}$, swarm size,
                              generations, seed and variable bounds \\
    Integration / output    & tolerances and maximum step
                              (Table~\ref{tab:numerics}), output sampling
                              interval, \texttt{SAVE\_PLOTS} and its target
                              directory \\
    \bottomrule
  \end{tabular}
\end{table}

Not every combination of switches is meaningful, and the selections are read in a definite
order rather than independently. Enabling the segmented mode overrides both the guidance-law
and the mission-architecture selections outright, since the schedule supplies the former and
the segmented solver fixes the latter; selecting the indirect law overrides the
mission-architecture selection alone, that architecture running its own dedicated search. Of
the settings that survive, the triangular kick profile is honored only under the
apogee-circularization architecture, the population-based solvers always applying the instantaneous
form (Section~\ref{ssec:dispatch}), and the direct architecture closes an orbit only for
those guidance laws capable of a single continuous insertion burn. Where such a dependency
would otherwise fail silently --- leaving the user with a plausible-looking trajectory
produced by settings that were never applied --- it is enforced explicitly, either by
raising on the incompatible combination or by warning and falling back to a defensible
default.

\subsection{Usage}
\label{ssec:usage}

A single mission is configured entirely in
\texttt{Input\_File/simulation\_parameters.py} and run through \texttt{main.py},
which selects the guidance law and mission architecture, runs the appropriate
optimizer, propagates the optimal trajectory, and writes the result arrays and
summary metrics. Vehicle and environment models are taken from the
\texttt{Auxiliary/} package, and no source-code edits are needed to change the
vehicle, mission, guidance law, or optimizer.

Every run writes a self-describing archive --- the trajectory, the scalar results and the full
configuration, including the version of the code --- so that it can be inspected, compared
with other runs or re-plotted later without being flown again. The cases of
Chapter~\ref{chapter:results_and_discussion} are defined and flown by
\texttt{run\_results\_matrix.py}, which collects their archives into a single table, and the
figures of that chapter are drawn from those archives by \texttt{Plots/results\_figures/}.
Two older batch runners (\texttt{all\_guidance\_plotting/} and
\texttt{guidance\_comparison/}) cover only the four classical laws and are superseded by it.
\end{comment}
\begin{comment}
\section{Guidance Laws as Implemented}
\label{sec:laws_implemented}

Chapter~\ref{chapter:guidance} presents each law as derived. This section states how the
laws are realized in the simulator where the derivation leaves a decision open.

Each law is a stateless function: given the state, the target and a time-to-go, it returns a
steering command, and whatever memory it needs is held by the simulation core and handed back
at its next call. This makes the laws interchangeable within the simulation and the comparison
between them fair. Appendix~\ref{chapter:appendix_guidance_options} describes the options not
used in this work: further time-to-go estimators, kick programs, and variants of several laws.

The gravity turn commands $\alpha = 0$ after the kick for the remainder of the flight, exactly
as derived.

The constant pitch rate law engages at second-stage ignition, with $\theta_{\mathrm{init}}$ equal
to the flight-path angle there, and the pitch rate $\dot\theta$ is its one free constant
(Table~\ref{tab:law_constants}). The ramp therefore shapes the second-stage arc rather than the
climb through the atmosphere. The apogee-circularization form of the law, presented in
Appendix~\ref{sec:app_cpr}, ramps from the vertical and so steers from the atmospheric arc
until insertion.

The constants $(a,b)$ of the exponential pitch law are searched by the optimizer
(Table~\ref{tab:law_constants}). Architectures that do not optimize them fix them by a shooting
solve (Appendix~\ref{sec:app_exp_shooting}).

\end{comment}

% =============================================================================
\section{Experimental Design}
\label{sec:experimental_design}

\subsection{One Factor at a Time}
\label{ssec:one_factor}

\subsection{The Cases}
\label{ssec:case_definitions}

\subsection{Reporting Conventions}
\label{ssec:reporting_template}

\begin{comment}
Two conventions apply to every number that follows. First, every swarm-based case is the
result of a single run at a fixed seed; no replication across seeds was performed, so small
differences between cases carry no estimate of their repeatability. Second, the velocity
target is the rotating-frame circular speed of Eq.~\eqref{eq:v_target}, which credits the
full eastward speed of the launch site as if the vehicle flew due east
(Chapter~\ref{ssec:frames}). The insertion state is therefore not a circular orbit but the apoapsis of an ellipse
whose periapsis lies below the target altitude, and an eccentricity computed with the same
convention reads zero by construction and cannot reveal it. The convention is shared by every
case, the reference included, so it does not favor one case over another; the $\Delta V$ budget
exposes it through its residual.
\end{comment}

%The tolerances, the step bound and the output sampling differ between execution paths. The values are given in Table~\ref{tab:numerics}, whose output grid is the sampling used for post-processing, not the internal step.
%The population-based paths use the tight setting because their objective scores the terminal radius directly, through the altitude penalty of Section~\ref{ssec:cost_function}: under the default tolerance a kilometer-level error in the radius is indistinguishable from a mis-targeted trajectory, and the swarm converges on the numerical artifact. Their maximum step is relaxed to \qty{10}{\s}, since the exo-atmospheric dynamics are smooth and the crash event remains reliably bracketed. The brute-force path keeps the default for its powered arcs: its cut-off is located by an event and a root search rather than scored by a penalty, so it is far less sensitive to the same error. 
%Its coast to apogee is the coast arc of the population-based paths and uses their settings (Section~\ref{ssec:grid_search}). 

\begin{comment}
    

\begin{table}[htbp]
  \centering
  \caption{Integration settings by execution path.}
  \label{tab:numerics}
  \begin{tabular}{l c c c l}
    \toprule
    \textbf{Execution path} & \textbf{rtol} & \textbf{atol} &
      \textbf{max.\ step} & \textbf{Output grid} \\
    \midrule
    Apogee circularization, powered arcs        & $10^{-3}$ (default) & $10^{-8}$ & \qty{1}{\s}  & \qty{0.01}{\s} \\
    Apogee circularization, coast to apogee     & $10^{-9}$ & $10^{-9}$ & \qty{10}{\s} & \qty{0.01}{\s} \\
    Coast / direct / indirect / segmented & $10^{-9}$ & $10^{-9}$ & \qty{10}{\s}
                                     & \qty{0.01}{\s} \\
    \bottomrule
  \end{tabular}
\end{table}
\end{comment}


